\usepackage{mathrsfs}
\let\mathcal\mathscr
\multlinegap0pt
\newcommand{\sbullet}{{\scriptscriptstyle\bullet}}
\newcommand{\csbullet}{{\raisebox{1pt}{$\sbullet$}}}
\newcommand{\Psfrac}[2]{(\sfrac{#1}{#2})}
\newcommand{\psfrac}[2]{\sfrac{(#1)}{#2}}
\newcommand{\spfrac}[2]{\sfrac{#1}{(#2)}}
\newcommand{\pspfrac}[2]{\sfrac{(#1)}{(#2)}}
\newcommand{\Ppsfrac}[2]{(\psfrac{#1}{#2})}
\newcommand{\llbracket}{[\![}
\newcommand{\rrbracket}{]\!]}
\newcommand{\Changel}{\mathcode`l="7160}
\newcommand{\Changelback}{\mathcode`l="716C}
\newcommand{\NoChangel}[1]{%
\expandafter\let\csname old\string#1\endcsname=#1
\let#1=\relax
\newcommand{#1}{\mathcode`l="716C\csname old\string#1\endcsname\mathcode`l="7160 }%
}

\NoChangel{\log}
\NoChangel{\ln}
\NoChangel{\lim}
\NoChangel{\limsup}
\NoChangel{\liminf}
\NoChangel{\varlimsup}
\NoChangel{\varliminf}
\NoChangel{\varinjlim}
\NoChangel{\varprojlim}

\usepackage{mathtools}
\usepackage{extarrows}

\newcommand{\C}{\mathbb{C}}

\theoremstyle{plain}
\newtheorem{theorem}{Theorem}[section]
\newtheorem{proposition}[theorem]{Proposition}
\newtheorem{corollary}[theorem]{Corollary}
\newtheorem{lemma}[theorem]{Lemma}
\theoremstyle{definition}
\newtheorem{definition}[theorem]{Definition}
\newtheorem{example}[theorem]{Example}
\newtheorem{notation}[theorem]{Notation}
\theoremstyle{remark}
\newtheorem{remark}[theorem]{Remark}

\newcommand{\Z}{\mathbb{Z}}
\newcommand{\cD}{\mathcal{D}}
\newcommand{\cS}{\mathcal{S}}
\newcommand{\cA}{\mathcal{A}}
\newcommand{\cE}{\mathcal{E}}
\newcommand{\cF}{\mathcal{F}}
\DeclareMathOperator{\Aut}{Aut}
\newcommand{\VEV}[1]{{\bigl\langle 0 \big| {#1} \big| 0 \bigr\rangle}}
\newcommand{\BVEV}[1]{{\Bigl\langle 0 \Big| {#1} \Big| 0 \Bigr\rangle}}
\newcommand{\VEVc}[1]{{\bigl\langle 0 \big| {#1} \big| 0 \bigr\rangle}^\circ}
\newcommand{\vac}{{\big| 0 \big\rangle}}
\newcommand{\res}{\mathop{\rm res}}
\newcommand{\hd}{\mathfrak{h}^\csbullet}
\newcommand{\hc}{\mathfrak{h}^\circ}
\newcommand{\Hd}{H^\csbullet}
\newcommand{\Hc}{H}
\newcommand{\cH}{\mathcal{H}}
\newcommand{\euro}{\mathbb{J}}
\DeclareMathOperator{\Id}{Id}
